\documentclass[11pt,a4paper]{article}

\usepackage[margin=1in]{geometry}
\usepackage{hyperref}
\usepackage{enumitem}
\usepackage{graphicx}
\usepackage{xcolor}
\usepackage{titling}

% -----------------------------------------------------
% ResQAI Project Proposal
% -----------------------------------------------------

\title{ResQAI: AI-Assisted Humanitarian Disaster Message Classification}

% Add subtitle
\pretitle{\begin{center}\bfseries\fontsize{18bp}{18bp}\selectfont}

\newcommand{\BvrSubtitle}[1]{%
  \posttitle{%
    \par\end{center}
    \begin{center}\large#1\end{center}
    \vskip 0.5em}%
}

% Hints in section titles
\newcommand{\BvrHint}[1]{%
  {\normalfont\normalsize\itshape\hfill #1}}

\BvrSubtitle{%
  Project Proposal \\
  Submitted to: \\[0.3cm]
  \bfseries Thapar Institute of Engineering and Technology
}

\author{
  Dhwani Agarwal%
  \thanks{Roll No: \texttt{1024031024}}
  \and
  Sanchita Jain%
  \thanks{Roll No: \texttt{1024031020}}
  \and
  Yuvakshi Sood%
  \thanks{Roll No: \texttt{1024031025}}
}

\date{\today}

\begin{document}

\maketitle

% =====================================================
% 1. HIGHER-ORDER GOAL
% =====================================================

\section{Higher-order goal%
  \BvrHint{\ldots secondary framing}}

This project aims to improve the efficiency of disaster information
processing by automatically identifying and organizing humanitarian
information from large volumes of social-media messages.

During natural disasters and humanitarian emergencies, affected individuals
often use social media to report injuries, request assistance, share
information about damaged infrastructure, communicate evacuation
requirements, report missing people, or offer rescue and volunteering
support.

The higher-order goal of ResQAI is to transform this unstructured and
high-volume information into structured, actionable categories that can
support emergency-response personnel and other decision makers.

The immediate engineering objective is to develop a deployable software
pipeline that can classify disaster-related messages rapidly and provide
additional contextual analysis for ambiguous messages.

% =====================================================
% 2. TIME-TO-VALUE
% =====================================================

\section{Time-to-value %
  \BvrHint{\ldots primary heuristic}}

\begin{itemize}[leftmargin=0.7cm]

    \item We will deliver meaningful increments quickly by first implementing
    the core user workflow and validating the classification component.

    \item We will prioritize a working end-to-end prototype over unnecessary
    infrastructure, allowing the core classification workflow to be
    demonstrated early.

    \item The first iteration will focus on message input, humanitarian
    classification, confidence estimation, and result presentation.

    \item Subsequent iterations will add contextual analysis, structured
    information extraction, and improved decision-support functionality.

    \item Success will be measured using classification accuracy, weighted
    F1 score, inference behavior, and the usefulness of the resulting
    structured information.

\end{itemize}

% =====================================================
% 3. PROBLEM STATEMENT
% =====================================================

\section{Problem Statement}

Natural disasters generate large amounts of social-media communication in a
short period of time. These messages may contain valuable information about
people, infrastructure, emergency needs, evacuation, rescue activities, and
safety conditions.

However, several challenges make it difficult to use this information
effectively:

\begin{itemize}[leftmargin=0.7cm]

    \item \textbf{Information overload:} A large number of messages may be
    generated during a disaster, making manual inspection difficult.

    \item \textbf{Unstructured information:} Messages are generally short,
    informal, and lack a standardized structure.

    \item \textbf{Mixed relevance:} Useful humanitarian information is mixed
    with political discussions, opinions, emotional reactions, news, and
    unrelated content.

    \item \textbf{Semantic ambiguity:} Several humanitarian categories can
    have overlapping meanings, making classification difficult.

    \item \textbf{Limited response time:} Emergency information may become
    less useful if it cannot be identified quickly.

\end{itemize}

\textbf{Why this matters now (contextual relevance):}
Social-media platforms can provide near-real-time information during
disasters, but the usefulness of this information depends on the ability to
filter, organize, and interpret it efficiently.

% =====================================================
% 4. PROPOSED SOLUTION
% =====================================================

\section{Proposed Solution %
\BvrHint{\ldots Software Proposition}}

\subsection{Overview}

ResQAI is proposed as an AI-assisted disaster information processing system
that automatically classifies incoming crisis-related messages into
humanitarian categories.

The system consists of the following major components:

\begin{itemize}[leftmargin=0.7cm]

    \item \textbf{Message input:} Accept a disaster-related social-media
    message.

    \item \textbf{BERTweet classifier:} Use a transformer model pretrained
    specifically for Twitter-style text and fine-tuned on humanitarian
    disaster messages.

    \item \textbf{Confidence estimation:} Produce the predicted category
    together with a confidence score.

    \item \textbf{LLM fallback:} Route ambiguous or low-confidence messages
    to a large language model such as Gemini for additional contextual
    analysis.

    \item \textbf{Structured output:} Present the resulting category and,
    where applicable, additional information such as urgency, location,
    affected people, and recommended action.

\end{itemize}

\subsection{Core workflow}

\begin{enumerate}[leftmargin=0.7cm]

    \item A disaster-related message is submitted to ResQAI.

    \item The message is tokenized using the BERTweet tokenizer.

    \item The fine-tuned BERTweet classifier predicts one of the ten
    humanitarian categories.

    \item The system obtains the prediction confidence.

    \item High-confidence predictions are accepted directly.

    \item Low-confidence predictions can be passed to the secondary Gemini
    analysis layer.

    \item The system presents the classification and relevant structured
    information to the user.

\end{enumerate}

\subsection{Humanitarian classification categories}

The implemented classifier uses ten categories:

\begin{itemize}[leftmargin=0.7cm]

    \item \texttt{caution\_and\_advice}

    \item \texttt{displaced\_people\_and\_evacuations}

    \item \texttt{infrastructure\_and\_utility\_damage}

    \item \texttt{injured\_or\_dead\_people}

    \item \texttt{missing\_or\_found\_people}

    \item \texttt{not\_humanitarian}

    \item \texttt{other\_relevant\_information}

    \item \texttt{requests\_or\_urgent\_needs}

    \item \texttt{rescue\_volunteering\_or\_donation\_effort}

    \item \texttt{sympathy\_and\_support}

\end{itemize}

% =====================================================
% 5. SOLUTION APPROACH
% =====================================================

\section{Solution Approach %
\BvrHint{\ldots Engineering focus}}

This is an engineering-oriented machine learning project focused on
developing a practical system for disaster information classification.

The project uses established machine learning and natural language processing
techniques rather than proposing a new transformer architecture.

\subsection{Dataset and preprocessing %
\BvrHint{\ldots data preparation}}

The project uses the HumAID humanitarian information dataset.

The dataset contains:

\begin{itemize}[leftmargin=0.7cm]

    \item 53,531 training messages,
    \item 7,793 validation messages,
    \item 15,160 test messages,
    \item 76,484 messages in total,
    \item 10 humanitarian classification categories.

\end{itemize}

Messages are tokenized using the BERTweet tokenizer.

The maximum sequence length is set to 128 tokens. Dataset analysis showed that
the average message length is approximately 35 tokens, while only a very
small proportion of messages reaches the 128-token limit.

\subsection{Model development %
\BvrHint{\ldots experimental engineering}}

Multiple approaches were evaluated before selecting the final model.

The investigated approaches included:

\begin{itemize}[leftmargin=0.7cm]

    \item TF-IDF with Naive Bayes,
    \item TF-IDF with Logistic Regression,
    \item class-weighted Logistic Regression,
    \item Linear SVM,
    \item Random Forest,
    \item DistilBERT,
    \item class-weighted DistilBERT,
    \item BERTweet.

\end{itemize}

BERTweet produced the strongest overall performance and was therefore selected
as the primary classifier.

\subsection{BERTweet fine-tuning %
\BvrHint{\ldots primary ML component}}

The final BERTweet configuration uses:

\begin{itemize}[leftmargin=0.7cm]

    \item Learning rate: $1\times10^{-5}$
    \item Maximum sequence length: 128
    \item Training batch size: 8
    \item Evaluation batch size: 8
    \item Maximum epochs: 3
    \item Weight decay: 0.01
    \item Warm-up steps: 500
    \item Linear learning-rate scheduler
    \item FP16 mixed precision during training
    \item Early stopping
    \item Best validation weighted F1 checkpoint selection

\end{itemize}

\subsection{Overfitting control %
\BvrHint{\ldots generalization}}

The model is evaluated after each epoch using validation data.

The best-performing checkpoint is selected using weighted F1 rather than
automatically using the final training epoch.

The experiments showed that validation performance peaked around the second
epoch and decreased slightly during the third epoch. The use of early
stopping and best-checkpoint selection therefore helps prevent unnecessary
additional training from reducing generalization.

\subsection{Confidence-based secondary analysis}

BERTweet is used as the primary classifier.

For messages where the model is sufficiently confident, the classification
result can be returned directly.

For ambiguous messages, a secondary LLM layer using Gemini is proposed.

The secondary layer can analyze:

\begin{itemize}[leftmargin=0.7cm]

    \item humanitarian category,
    \item urgency,
    \item location,
    \item number of people affected,
    \item requested assistance,
    \item recommended action.

\end{itemize}

The LLM therefore complements rather than replaces the trained
domain-specific classifier.

% =====================================================
% 6. SYSTEM ARCHITECTURE
% =====================================================

\section{System Architecture %
\BvrHint{\ldots software-based solution}}

The proposed ResQAI architecture consists of several logical components.

\begin{itemize}[leftmargin=0.7cm]

    \item \textbf{Input layer:} Receives a social-media or crisis message.

    \item \textbf{NLP layer:} Tokenizes and preprocesses the incoming
    message.

    \item \textbf{BERTweet layer:} Performs humanitarian classification.

    \item \textbf{Confidence layer:} Determines whether the classifier is
    sufficiently confident in its prediction.

    \item \textbf{LLM layer:} Provides additional analysis for ambiguous
    messages.

    \item \textbf{Output layer:} Displays classification, confidence, and
    additional structured information.

\end{itemize}

The conceptual architecture is:

\begin{verbatim}
                    Crisis Message
                          |
                          v
                 +----------------+
                 | Preprocessing  |
                 |  + Tokenizer   |
                 +-------+--------+
                         |
                         v
                 +----------------+
                 |    BERTweet    |
                 |   Classifier   |
                 +-------+--------+
                         |
                  Category +
                  Confidence
                         |
              +----------+----------+
              |                     |
       High Confidence        Low Confidence
              |                     |
              v                     v
       Direct Result          Gemini API
                                    |
                                    v
                            Contextual Analysis
                                    |
                                    v
                           Structured Output
\end{verbatim}

% =====================================================
% 7. EVALUATION CRITERION
% =====================================================

\section{Evaluation Criterion %
\BvrHint{\ldots measurable and attributable}}

The success of the project will be evaluated using measurable machine
learning and system-level metrics.

\subsection{Primary evaluation metrics}

\textbf{Classification Accuracy:}

\begin{equation}
Accuracy =
\frac{\text{Correct Predictions}}
{\text{Total Predictions}}
\end{equation}

Accuracy measures the overall proportion of correctly classified messages.

\textbf{Weighted F1 Score:}

The F1 score combines precision and recall.

\begin{equation}
F1 = 2 \times
\frac{Precision \times Recall}
{Precision + Recall}
\end{equation}

The weighted average accounts for the different support of individual
classes.

\subsection{Current achieved performance}

The final BERTweet model achieved:

\begin{itemize}[leftmargin=0.7cm]

    \item \textbf{Best Validation Accuracy: 78.40\%}
    \item \textbf{Best Validation Weighted F1: approximately 78.05\%}
    \item \textbf{Test Accuracy: 78.67\%}
    \item \textbf{Test Weighted F1: 78.27\%}

\end{itemize}

The test results were obtained from the unseen test set containing 15,160
messages.

\subsection{Secondary evaluation metrics}

The following metrics can be considered during further system development:

\begin{itemize}[leftmargin=0.7cm]

    \item \textbf{Per-class precision and recall:} To identify categories
    that require improvement.

    \item \textbf{Confusion analysis:} To identify commonly confused
    categories.

    \item \textbf{Inference time:} Time required to classify a single
    message.

    \item \textbf{LLM fallback rate:} Percentage of messages requiring
    secondary analysis.

    \item \textbf{Confidence calibration:} Relationship between model
    confidence and actual prediction correctness.

\end{itemize}

\subsection{Validation plan %
\BvrHint{\ldots high-level}}

The model will be evaluated using the predefined training, validation, and
test splits.

The validation set is used for model selection and hyperparameter decisions,
while the test set is reserved for final performance evaluation.

This separation prevents the test set from influencing model selection.

% =====================================================
% 8. SCALABILITY
% =====================================================

\section{Scalability %
\BvrHint{\ldots with foresight}}

The proposition should scale in both theoretical and practical senses.

\begin{enumerate}[leftmargin=0.7cm]

    \item \textbf{Scalable classification:}

    The trained classifier can process incoming messages independently,
    allowing multiple inference requests to be handled without retraining the
    model.

    \item \textbf{Batch processing:}

    Messages can be grouped into batches for more efficient processing when
    large volumes of information are available.

    \item \textbf{Separation of components:}

    The classifier, LLM analysis layer, and user interface can be separated
    into independent components.

    \item \textbf{API-based architecture:}

    The classification model can be exposed through an API so that different
    interfaces or applications can use the same inference service.

    \item \textbf{LLM cost control:}

    A confidence-based fallback mechanism avoids sending every message to an
    external LLM, which can reduce unnecessary API usage.

\end{enumerate}

% =====================================================
% 9. ENGINE AVAILABILITY
% =====================================================

\section{Engine availability heuristic}

A mature set of machine learning and software components is available to
implement the proposed system:

\begin{itemize}[leftmargin=0.7cm]

    \item \textbf{Hugging Face Transformers:} Model loading, tokenization,
    and transformer inference.

    \item \textbf{PyTorch:} Model execution and machine learning operations.

    \item \textbf{BERTweet:} Primary humanitarian text-classification model.

    \item \textbf{Gemini API:} Proposed secondary contextual reasoning layer.

    \item \textbf{Python:} Core implementation and integration language.

    \item \textbf{Web/API framework:} Can be used to expose the model as an
    application service.

    \item \textbf{Cloud or local deployment:} The inference component can be
    deployed on suitable compute infrastructure depending on the required
    scale.

\end{itemize}

The project uses established tools and models so that the engineering effort
can focus on system integration, model evaluation, and practical usability.

% =====================================================
% 10. PROJECT SCOPE AND DELIVERABLES
% =====================================================

\section{Project scope and deliverables %
\BvrHint{\ldots iteration-friendly}}

\subsection{Initial deliverable %
\BvrHint{\ldots first iteration}}

The initial implementation includes:

\begin{itemize}[leftmargin=0.7cm]

    \item HumAID dataset preparation.

    \item Training and evaluation of multiple baseline classifiers.

    \item Fine-tuning of BERTweet for ten-class humanitarian classification.

    \item Validation and test evaluation.

    \item Prediction of humanitarian category for a new message.

    \item Confidence score generation.

    \item Model checkpoint preservation for inference.

\end{itemize}

\subsection{Subsequent deliverables}

Further iterations will include:

\begin{itemize}[leftmargin=0.7cm]

    \item User-facing message input interface.

    \item Confidence-based Gemini fallback.

    \item Structured output generation.

    \item Urgency and priority identification.

    \item Location extraction.

    \item Recommended-action generation.

    \item Visualization of classification results.

    \item Optional storage of analyzed messages and results.

\end{itemize}

% =====================================================
% 11. RISKS AND MITIGATIONS
% =====================================================

\section{Risks and mitigations}

\begin{itemize}[leftmargin=0.7cm]

    \item \textbf{Classification errors:}

    The classifier may incorrectly categorize ambiguous messages.

    \textbf{Mitigation:} Use confidence scores, detailed validation, confusion
    analysis, and secondary LLM analysis for low-confidence cases.

    \item \textbf{Semantic overlap between classes:}

    Some humanitarian categories have inherently similar meanings.

    \textbf{Mitigation:} Analyze confusion patterns and consider hierarchical
    or multilabel classification in future versions.

    \item \textbf{LLM dependency:}

    The Gemini fallback requires external API access.

    \textbf{Mitigation:} Keep BERTweet as the primary classifier and use the
    LLM only when required.

    \item \textbf{Incorrect emergency interpretation:}

    Automated predictions should not directly trigger critical emergency
    actions without verification.

    \textbf{Mitigation:} Treat the system as decision support and maintain
    human oversight for high-impact decisions.

    \item \textbf{Data privacy:}

    Social-media messages may contain sensitive information.

    \textbf{Mitigation:} Minimize stored information, avoid unnecessary
    personally identifiable information, and apply appropriate access
    controls.

    \item \textbf{Computational requirements:}

    Transformer inference requires more resources than simple TF-IDF
    classifiers.

    \textbf{Mitigation:} Use batching, model optimization, and appropriate
    deployment infrastructure when scaling.

\end{itemize}

% =====================================================
% 12. EXPECTED OUTCOME
% =====================================================

\section{Expected outcome}

The expected outcome of the project is a working AI-assisted prototype that
can accept a disaster-related message and automatically determine its most
likely humanitarian category.

The primary machine learning component has already demonstrated a test
accuracy of approximately 78.67\% and a weighted F1 score of approximately
78.27\%.

The completed system is expected to provide:

\begin{itemize}[leftmargin=0.7cm]

    \item automated humanitarian classification,
    \item confidence-based prediction,
    \item identification of ambiguous messages,
    \item optional LLM-based contextual analysis,
    \item structured disaster information,
    \item and a foundation for future real-time disaster information
    processing.

\end{itemize}

% =====================================================
% 13. SUMMARY
% =====================================================

\section{Summary}

ResQAI proposes an AI-assisted approach for organizing and interpreting
large volumes of disaster-related social-media messages.

The project focuses on time-to-value by first developing and validating the
core classification model before adding more advanced system components.

The primary model, BERTweet, was fine-tuned using the HumAID dataset and
achieved 78.67\% accuracy and 78.27\% weighted F1 on the unseen test set.

The system is designed around a practical hybrid architecture: BERTweet
provides fast domain-specific classification, while a large language model
such as Gemini can be used for additional contextual reasoning when the
classifier is uncertain.

The project therefore combines machine learning, natural language processing,
and software-system integration to create a practical information-processing
tool for disaster-response scenarios.

The final objective is not to replace emergency responders but to reduce the
amount of unstructured information they need to manually process and help
surface potentially relevant humanitarian information more efficiently.

\end{document}